\paragraph{Robustness sweeps.}\label{sec:sweeps}
Table~\ref{tab:sweeps} varies one factor at a time on Fashion-MNIST: the strength of the concept shift, the number of groups, the fraction of participating clients and the number of clients. Baselines are given the true $K$ where they need it.

\begin{tableorg}[!htbp]
\centering
\caption{Robustness sweeps (Fashion-MNIST, 3 seeds). For DisCo-CFL: ARI, number of clusters, conflict-merge rate, accuracy and concept accuracy. For comparison: the best baseline on accuracy and on concept accuracy (in parentheses), and the Oracle's accuracy.}\label{tab:sweeps}
\scriptsize
\setlength{\tabcolsep}{3pt}
\resizebox{\textwidth}{!}{%
\begin{tabular}{llcccccccc}
\toprule
& & \multicolumn{5}{c}{DisCo-CFL} & & & \\
\cmidrule(lr){3-7}
Factor & Setting & ARI & $K$ & Conflict\% & Acc & Concept & Best baseline Acc & Best baseline Concept & Oracle Acc \\
\midrule
Strength & $\theta=15^\circ$ & 0.42 & 3.0 & 32.1 & 90.0 & 82.8 & 90.2 (FeSEM) & 83.0 (FeSEM) & 90.5 \\
Strength & $\theta=30^\circ$ & 0.93 & 4.3 & 1.7 & 90.3 & 83.1 & 90.0 (PACFL) & 81.2 (PACFL) & 90.4 \\
Strength & $\theta=45^\circ$ & 1.00 & 4.0 & 0.0 & 90.4 & 83.8 & 89.9 (FeSEM) & 81.8 (FeSEM) & 90.4 \\
Strength & $\theta=90^\circ$ & 1.00 & 4.0 & 0.0 & 90.5 & 83.9 & 90.2 (MTCFL) & 80.6 (FeSEM) & 90.5 \\
Groups & $K=2$ & 0.94 & 3.3 & 0.0 & 91.4 & 86.5 & 88.5 (IFCA) & 80.2 (IFCA) & 91.5 \\
Groups & $K=4$ & 1.00 & 4.0 & 0.0 & 90.8 & 84.9 & 90.6 (FedAvg-FT) & 80.9 (FedAvg-FT) & 90.8 \\
Groups & $K=8$ & 0.92 & 8.0 & 0.3 & 90.4 & 81.7 & 88.0 (FeSEM) & 75.4 (FeSEM) & 90.5 \\
Participation 30\% & rotation & 1.00 & 4.0 & 0.0 & 89.1 & 78.4 & 87.5 (PACFL) & 75.1 (FeSEM) & 89.1 \\
Participation 30\% & label swap & 0.98 & 4.0 & 0.0 & 89.3 & 78.3 & 86.5 (FeSEM) & 73.0 (FeSEM) & 89.3 \\
Scale & $N=200$ & 0.98 & 6.3 & 0.0 & 91.1 & 86.6 & 89.8 (PACFL) & 83.9 (FeSEM) & 91.1 \\
\bottomrule
\end{tabular}}

\end{tableorg}

\paragraph{Strength of the shift.} With rotations of $0,\theta,2\theta,3\theta$ degrees, DisCo-CFL separates the groups perfectly from $\theta=45^\circ$ on and almost perfectly at $30^\circ$ (ARI $0.93$). At $15^\circ$ neighbouring concepts are so close that their class-conditional gradients agree above the threshold, and DisCo-CFL merges them (ARI $0.42$). This is the behaviour predicted by the separation assumption: the method groups what it cannot tell apart. It costs little, because nearby concepts can share a model: accuracy is 90.0\% against 90.5\% for the Oracle, and 0.2 points below FeSEM, which is given $K=4$. For every $\theta\ge30^\circ$, DisCo-CFL has the best accuracy and concept accuracy and is within $0.1$ points of the Oracle.

\paragraph{Number of groups.} With $K=2$, $4$ and $8$ label-permutation groups, DisCo-CFL finds $K$ without being told (3.3, 4.0 and 8.0 clusters; the extra cluster at $K=2$ is a conflict-free split) and merges at most 0.3\% of conflicting pairs. Its accuracy is within $0.1$ points of the Oracle for all $K$. The margin over the best baseline is 2.9, 0.2 and 2.4 points in accuracy and 6.3, 4.0 and 6.3 points in concept accuracy for $K=2,4,8$. With $K=2$ and $K=8$ the baselines given $K$ confuse groups (IFCA and FeSEM merge 25\% and 4\% of the conflicting pairs), whereas with $K=4$ the fine-tuned global model is a strong competitor on local accuracy only.

\paragraph{Partial participation.} When only 30\% of the clients take part in each round, the warm-up model is noisier. DisCo-CFL still recovers the groups (ARI $1.00$ and $0.98$) and matches the Oracle, while the best baseline is 1.6--2.8 points below in accuracy and 3.3--5.3 points below in concept accuracy.

\paragraph{Number of clients.} With 200 clients of 200 samples each, DisCo-CFL reaches ARI $0.98$ with no conflicting merge. The extra clusters ($K=5$--$7$) are conflict-free splits of groups; three of them are singletons flagged as anomalous in the audit log. Accuracy equals the Oracle's (91.1\%), 1.3 points above the best baseline, and concept accuracy is 2.7 points above the best baseline. This is in line with the clustering-only scalability study of Section~\ref{sec:cost}.

\paragraph{Statistical significance.}
Table~\ref{tab:signif} compares DisCo-CFL with each baseline over all (scenario, seed) pairs with concept shift that both methods were run on. DisCo-CFL is significantly better than every baseline on mean accuracy, concept accuracy and worst-group accuracy (one-sided Wilcoxon signed-rank tests, all $p<0.01$). The closest competitors are FeSEM, MTCFL and FedAvg-FT. MTCFL and FedAvg-FT win on local accuracy mainly on CIFAR-10 (5--6 of the 6 CIFAR-10 pairs), and FeSEM wins on 8 of 57 pairs, mostly CIFAR-10 rotation and the $15^\circ$ rotation. On concept accuracy DisCo-CFL wins against every baseline in at least 84\% of the pairs.

\begin{table}[!htbp]
\centering
\caption{Paired one-sided Wilcoxon signed-rank tests of DisCo-CFL against each baseline over all (scenario, seed) pairs with concept shift. W/T/L: wins, ties ($|\Delta|<0.05$ points) and losses of DisCo-CFL.}\label{tab:signif}
\small
\begin{tabular}{lccccccc}
\toprule
& & \multicolumn{2}{c}{Accuracy} & \multicolumn{2}{c}{Concept accuracy} & \multicolumn{2}{c}{Worst-group accuracy} \\
\cmidrule(lr){3-4}\cmidrule(lr){5-6}\cmidrule(lr){7-8}
Baseline & pairs & W/T/L & $p$ & W/T/L & $p$ & W/T/L & $p$ \\
\midrule
FedAvg & 57 & 57/0/0 & 2.6e-11 & 57/0/0 & 2.6e-11 & 57/0/0 & 2.6e-11 \\
Local & 33 & 33/0/0 & 1.2e-10 & 33/0/0 & 1.2e-10 & 32/0/1 & 2.3e-10 \\
FedAvg-FT & 33 & 27/1/5 & 2.3e-5 & 33/0/0 & 1.2e-10 & 24/0/9 & 0.004 \\
IFCA & 54 & 54/0/0 & 8.1e-11 & 54/0/0 & 8.1e-11 & 54/0/0 & 8.1e-11 \\
FeSEM & 57 & 45/4/8 & 2.5e-8 & 48/2/7 & 1.4e-8 & 47/3/7 & 5.2e-8 \\
MTCFL & 39 & 30/2/7 & 0.002 & 39/0/0 & 1.8e-12 & 30/0/9 & 1.4e-4 \\
FL+HC & 57 & 57/0/0 & 2.6e-11 & 57/0/0 & 2.6e-11 & 57/0/0 & 2.6e-11 \\
PACFL & 57 & 54/0/3 & 7.0e-11 & 57/0/0 & 2.6e-11 & 53/1/3 & 5.1e-11 \\
\bottomrule
\end{tabular}

\end{table}
